\begin{abstract}
  Lack of motivation and disinterest are the main obstacles to consistency in physical exercise. To address this problem, this study proposes a mobile app integrated with smartwatches, focused on the gamification of physical activities. \\
  \indent The system uses real-time biometric data monitoring (biofeedback) to drive an immersive survival narrative, transforming physical effort into digital progress. Based on Flow Theory and Dynamic Difficulty Adjustment (DDA), the application structures gamified challenges and uses real-time biometric data to support activity monitoring. This synergy prevents disinterest due to ease or frustration due to excessive effort. \\
  \indent Anchored in Self-Determination Theory (SDT), the integration of immersion and biometric precision seeks to convert extrinsic motivation into intrinsic motivation. The ultimate goal is to promote greater engagement, autonomy, and sustainable adherence to exercise.
\end{abstract}

\begin{keywords}
Gamification, Biofeedback, Exergames, Dynamic Difficulty Adjustment.
\end{keywords}

\begin{resumo}
   A falta de motivação e o desinteresse são os principais obstáculos à constância nos exercícios físicos. Para amenizar esse problema, este trabalho propõe um aplicativo móvel integrado a \textit{smartwatches}, focado na gamificação de atividades físicas. \\
   \indent O sistema utiliza o monitoramento de dados biométricos em tempo real (\textit{biofeedback}) para conduzir uma narrativa de sobrevivência imersiva, transformando o esforço físico em progresso digital. Fundamentada na Teoria do Fluxo e no Ajuste Dinâmico de Dificuldade (DDA), a aplicação estrutura desafios gamificados e utiliza dados biométricos em tempo real como suporte ao acompanhamento da atividade. Essa sinergia evita o desinteresse por facilidade ou a frustração por excesso de esforço. \\
   \indent Ancorada na Teoria da Autodeterminação (SDT), a integração entre imersão e precisão biométrica busca converter a motivação extrínseca em intrínseca. O objetivo final é promover maior engajamento, autonomia e uma adesão sustentável à prática de exercícios. 
\end{resumo}

\begin{palavrasChaves}
Gamificação, Biofeedback, Exergames, Ajuste Dinâmico de Dificuldade.
\end{palavrasChaves}

